\section{Christ in the Person Before Us}

In Matthew~25, Christ judges our response to people whom we can serve now. The Son of Man comes in glory, gathers the nations, and speaks of hunger, thirst, nakedness, estrangement, sickness, and imprisonment. He identifies himself with the least of his brethren. What was done or refused to them concerns him (Mt 25:31--46). His identification makes the refusal of mercy a refusal of the King before whom all stand. The passage asks us to recognize his claim in a person's need before the judgment makes it manifest.

The corporal works of mercy require practical attention to a person's needs. Food must be obtained and prepared; a visit requires time; welcome can require making room in a household. Such care may need to continue after its first urgency has passed. Our willingness can weaken when a need becomes familiar or interrupts plans we value. Christ's words give a reason to persevere: he receives what is done for the least as done for himself.

Chrysostom asks his hearers to notice how accessible the commanded works are. Christ does not say that they cured the sick or broke open the prison. He says that they visited and came. The hungry one did not demand a costly banquet, but necessary food (LXXIX.1--2). The distinction frees charity from an excuse familiar in every age. We can decline a possible service because we cannot accomplish a complete rescue. Yet the inability to remove every suffering does not erase the duty to do the good within reach.

For example, a household may be able to help a family through one difficult week while lacking the means to provide for it indefinitely. The limited duration should be made clear, so that an offer of help does not become a promise which cannot be kept. Further assistance may require cooperation with others. These limits call for prudence, but the good which can responsibly be done remains possible before a lasting remedy has been found.

\ifdefined\TriptychPrintEdition
\looseness=1
\else
\fi
The condemnations in the passage include omission. The refusal to give food, clothing, welcome, or a visit is enough to show a terrible estrangement from Christ. Chrysostom joins this to the buried talent and the unprepared virgins: neglect of the neighbor can ruin a life which congratulates itself for avoiding more conspicuous wrongs (LXXVIII.1--3). We therefore examine more than the injuries we deliberately inflicted. We also examine the duties we evaded and the needs we freely chose to ignore.\par

This examination must remain concrete. The neighbor we have neglected may be a relative whose illness has become tiresome, a person whose need interrupts our preferred work, or someone from whom we expect no gratitude. We may have made generosity dependent upon being noticed. We may have confined it to people who share our tastes and opinions. We may support distant good causes while using impatience to keep a dependent person from approaching us. Christ's identification with the least reaches these arrangements and asks what kind of love they contain.

James begins his account of living faith by rebuking partiality toward the well-dressed rich man and contempt toward the poor man (Jas 2:1--13). The person who can benefit us receives a place of honor, while the inconvenient person is made to stand aside. Such preference can survive under courteous language and respectable habits. An annual examination before judgment can uncover the social calculations by which we decide whose suffering deserves our attention. The Lord's presence is not proportioned to a person's usefulness to us.

Chrysostom's own exposition of Christ's least brethren emphasizes the poor and neglected believer, including the lay Christian and not only the monk (LXXIX.1). The claim gives the baptized neighbor a searching place within the Church's life. Baptism does not make another member's deprivation someone else's problem. At the same time, Christ commands love of enemies and shows mercy beyond the boundaries by which we ordinarily recognize friends (Mt 5:43--48; Lk 10:25--37). Our charity must be as wide as the Lord's command, beginning with the actual persons he places before us.

Merciful assistance requires listening. We may discover that the help we enjoy offering differs from the help needed. Prudence governs the form of a gift, responsibilities to dependents remain real, and some requests cannot rightly be fulfilled. When a request must be refused, care for the person can continue through seeking another form of help. We should examine whether we have made that effort or used a real limitation as a reason to avoid further inconvenience.

The good done in secret belongs here as fully as the mercy made visible. A person can serve for years without receiving the affection hoped for. The Lord does not forget the repeated visit, the patient washing, the money quietly given, or the interrupted sleep. Cyril's assurance that no faithful act is lost before God guards these works from bitterness (XV.23). We may desire natural gratitude, but we need not make gratitude a condition of obedience. Christ will know the whole service, including the effort by which love endured its obscurity.

The judgment scene ends in eternal punishment and eternal life. Christ has made our treatment of the neighbor a place of response to himself. Meditation should therefore lead to some care actually given. Where neglect has become habitual, one decision may need to include how that care will continue: when to visit, how much assistance can be offered, or who can help fulfill the obligation. The person's need, rather than the satisfaction of making a generous resolution, should determine what we undertake.

